% Modelo de relatório técnico HRAI-GO (ABNT NBR 10719:2015).
% Copie este arquivo para docs/relatorios/<tema>/relatorio.tex e preencha os campos marcados.
% Compilação: scripts/build-relatorios.sh <arquivo.tex>   (XeLaTeX + BibTeX via latexmk)
% Regras: docs/PADRAO_RELATORIOS.md
\documentclass[12pt,a4paper,oneside,english,brazil]{abntex2}
\usepackage{hraigo-abnt}

% ------------------------------------------------------------ Identificação (NBR 10719, 4.2.1.1.1)
\titulo{Título do relatório}
\subtitulo{subtítulo, se houver}
\codigo{RT-XXX-00}                       % código na série HRAI-GO (ver PADRAO_RELATORIOS.md)
\versao{1.0}
\data{2026}                               % ano de publicação (algarismos arábicos)
% \autor{...} \qualificacao{...} \contato{...} % padrão: Jefferson Wellington da Cunha, Doutorando
% \patrocinador{...}                       % padrão: AeonBridge Co. (https://aeonbridge.com/)
% \semagradecimentospatrocinio             % omite os Agradecimentos com as bases do patrocínio
% \local{Cidade}                           % padrão: Brasil
% \classificacao{Ostensivo (público)}      % padrão
% \equipe{Nome -- função \\ Nome -- função} % opcional: imprime o verso da folha de rosto
\palavraschave{Governança de IA. Inteligência regulatória. Ontologia.}

\begin{document}
\selectlanguage{brazil}

% ============================================================ ELEMENTOS PRÉ-TEXTUAIS
\pretextual
\imprimirfolhaderostorelatorio            % obrigatório

\begin{resumorelatorio}                   % obrigatório (NBR 6028): objetivo, método, resultados, conclusões
Texto do resumo em parágrafo único, com 150 a 500 palavras, na terceira pessoa e na voz ativa.
Ele indica o objetivo, o método, os resultados e as conclusões do relatório.
\end{resumorelatorio}

% Listas opcionais: use quando houver mais de poucas ilustrações, tabelas ou siglas.
\pdfbookmark[0]{\listfigurename}{lof}
\listoffigures*
\cleardoublepage
\pdfbookmark[0]{\listtablename}{lot}
\listoftables*
\cleardoublepage

\begin{siglas}
  \item[ABNT] Associação Brasileira de Normas Técnicas
  \item[HRAI-GO] Human-centred Regulatory AI, Governance Observatory
\end{siglas}

\pdfbookmark[0]{\contentsname}{toc}
\tableofcontents*                          % obrigatório
\cleardoublepage

% ============================================================ ELEMENTOS TEXTUAIS
\textual

\chapter{Introdução}
Apresente o objetivo, o escopo e a justificativa do relatório. Cite as fontes no sistema
autor-data \cite{nbr10719} ou com o autor no texto, conforme \citeonline{pbia2025}.

\chapter{Desenvolvimento}
\section{Método}
Descreva o método de forma que outra equipe possa reproduzi-lo.

\section{Resultados}
\begin{table}[htb]
  \caption{Título da tabela, acima dela}\label{tab:exemplo}
  \centering
  \begin{tabular}{lr}
    \toprule
    Item & Valor\\
    \midrule
    Exemplo & 1\\
    \bottomrule
  \end{tabular}
  \fonte{Elaborado pelos autores (2026).}
\end{table}

A Tabela~\ref{tab:exemplo} mostra o formato de tabela: legenda acima, fonte abaixo, sem linhas verticais.

\begin{figure}[htb]
  \caption{Título da figura, acima dela}\label{fig:exemplo}
  \centering
  \begin{tikzpicture}
    \node[draw=lagoa,rounded corners,minimum width=4cm,minimum height=1cm] {Exemplo};
  \end{tikzpicture}
  \fonte{Elaborado pelos autores (2026).}
\end{figure}

\chapter{Considerações finais}
Retome os objetivos e apresente as conclusões e as recomendações.

% ============================================================ ELEMENTOS PÓS-TEXTUAIS
\postextual
\bibliography{referencias}                 % obrigatório

\begin{apendicesenv}                       % opcional: material elaborado pelos autores
\partapendices
\chapter{Exemplo de apêndice}
Conteúdo complementar elaborado pelos autores.
\end{apendicesenv}

% Anexos (opcional): material não elaborado pelos autores, em \begin{anexosenv}...\end{anexosenv}

\imprimirformularioidentificacao[Resumo do relatório para o formulário.]  % opcional, último elemento
\end{document}
